\section{The original Laplacian control along the exact period deformation}
\label{sec:period-laplacian-control-bridge}

We calculate the control parameter for the actual companion deformation,
and then transport its complete numerical range through the periods.
All metrics and coefficient vectors below are the original ones.

Fix the original cyclic algebra, admitted degree and positive Gram
\[
 E=\mathbb C[S]/(\chi),\quad q=\deg\chi\ge1,\quad
 A=M_S,\quad G=G_N>0,\quad k\in\mathbb R,
 \quad e_0=[1],\quad
 \ell[P]=[S^{q-1}]\operatorname{rem}_{\chi}P.
 \tag{LC1}
\]
Put \(R=e_0\ell\), \(A_t=A+tR\), and retain the actual endomorphisms
\[
 W_t=A_t^*G+GA_t-kG,\quad X_t=G^{-1}W_t,\quad
 L_t=A_t^2-kA_t,\quad B_t=A_t-\tfrac k2I.
 \tag{LC2}
\]
Here \(X_t\) is self-adjoint for
\(\langle v,w\rangle_G=v^*Gw\). Its finite operator norm is defined
by the supremum over the original nonzero vectors. Write
\(\epsilon_N=\|X_0\|_G\). This is an actual finite quantity, so
\(-\epsilon_NG\preceq W_0\preceq\epsilon_NG\) follows from the
spectral theorem for the displayed positive metric.

For the original reflection-stable source, its complete rank-two
calculation (TVB.22)--(TVB.30) evaluates this same quantity. In the
unchanged source-orthogonal columns \(b_j\) and norms \(\omega_j\),
put
\[
 a=b_N^*Gb_N,\quad d_+=b_{N+1}^*Gb_{N+1},\quad
 z_N=b_N^*Gb_{N+1},\quad D_N=\det K_N.
\]
The retained identities give
\[
 \epsilon_N^2=\frac{ad_+-|z_N|^2}{\omega_N^2}
       =\alpha_{N+1}\frac{\Delta_N}{D_N},
 \quad\alpha_{N+1}=\frac{\omega_{N+1}}{\omega_N},
 \quad
 \Delta_N=D_{N+1}+D_{N-1}-D_N-\widetilde D_N.
 \tag{LC3}
\]
The omitted-column determinant \(\widetilde D_N\), every boundary
case and its positive source-minor sum are defined and proved in
that full chapter. To check the radius here directly, the original
control is
\(W_0=G(b_{N+1}b_N^*+b_Nb_{N+1}^*)G/\omega_N\).
Its relative self-adjoint operator has rank at most two and trace
zero; \(z_N\) is purely imaginary by the retained reflection phase.
Expansion of its squared trace gives
\(2(ad_+-|z_N|^2)/\omega_N^2\), so its two possible nonzero
eigenvalues are exactly \(\epsilon_N,-\epsilon_N\).
The two-column determinant identity proves the second equality
in (LC3). No estimate or limiting assumption is needed to evaluate it.

\subsection{The exact additional radius of the companion term}

The adjoint of \(R\) for the original metric is
\(R^{\dagger_G}=G^{-1}\ell^*e_0^*G\). Thus
\[
 X_t-X_0=tR+\bar tR^{\dagger_G}.
 \tag{LC4}
\]
For \(q>1\), \(\ell e_0=0\), hence \(R^2=0\). Put
\[
 f=G^{-1}\ell^*,\quad a_0=e_0^*Ge_0>0,\quad
 b_0=\ell G^{-1}\ell^*>0,\quad
 \sigma_N=\sqrt{a_0b_0}.
 \tag{LC5}
\]
The columns \(e_0,f\) are \(G\)-orthogonal because
\(e_0^*Gf=\overline{\ell e_0}=0\). They are independent and their
Gram is exactly \(\operatorname{diag}(a_0,b_0)\). In these original
unscaled columns the matrix of (LC4) is
\[
 \begin{pmatrix}0&t b_0\\\bar t a_0&0\end{pmatrix}.
 \tag{LC6}
\]
Indeed \(Rf=b_0e_0\), \(R e_0=0\),
\(R^{\dagger_G}e_0=a_0 f\) and
\(R^{\dagger_G}f=0\). On their \(G\)-orthogonal complement both
terms vanish. The characteristic polynomial on the displayed
two-dimensional space is \(x^2-|t|^2a_0b_0\). Therefore the exact
additional operator norm is \(|t|\sigma_N\); all remaining
eigenvalues are zero. This retains the two original vector lengths,
with no choice of unit vectors.

For \(q=1\), instead \(R=I_E\), so (LC4) equals
\(2\operatorname{Re}(t)I_E\) and its exact norm is
\(2|\operatorname{Re}t|\). Consequently the entirely specified
quantity
\[
 E_N(t)=\epsilon_N+
 \begin{cases}|t|\sigma_N,&q>1,\\
               2|\operatorname{Re}t|,&q=1
 \end{cases}
 \quad\text{satisfies}\quad
 \|X_t\|_G\le E_N(t),\qquad
 -E_N(t)G\preceq W_t\preceq E_N(t)G.
 \tag{LC7}
\]
The inequality is the triangle inequality for the same operator
norm. It is not an assertion that equality always holds. The two
summands in (LC7) are the actual norms just computed, and the
original \(\epsilon_N\) remains available through (LC3).
The full Laplacian change is also exact:
\[
 L_t-L_0=t(AR+RA-kR)+t^2R^2.
 \tag{LC8}
\]
Thus the quadratic term vanishes for \(q>1\), while for \(q=1\)
it is the retained scalar \(t^2I_E\).

\subsection{An explicit numerical-range enclosure with no missing input}

Expansion of (LC2), without a commutation assumption, proves
\[
 L_t^*G-GL_t=A_t^*W_t-W_tA_t,\qquad
 GL_t=W_tB_t-B_t^*GB_t-\tfrac{k^2}{4}G.
 \tag{LC9}
\]
For the second equality insert
\(W_t=B_t^*G+GB_t\), so the right side is
\(GB_t^2-k^2G/4=GL_t\). The first equality follows by expansion;
the two middle terms \(A_t^*GA_t\) cancel.

For each unchanged nonzero \(v\in E\), define the homogeneous
observations
\[
 z=\frac{v^*GL_tv}{v^*Gv},\quad
 r=\frac{\|B_tv\|_G}{\|v\|_G},\quad
 a+ib=\frac{\langle v,X_tB_tv\rangle_G}{\|v\|_G^2}.
 \tag{LC10}
\]
Cauchy--Schwarz and (LC7) give
\(a^2+b^2\le E_N(t)^2r^2\). The second identity in (LC9) gives
\(\operatorname{Re}z=a-r^2-k^2/4\) and
\(\operatorname{Im}z=b\). Write \(e=E_N(t)\). First,
\(a\le er\) implies
\(\operatorname{Re}z\le(e^2-k^2)/4-(r-e/2)^2\).
Second, direct multiplication gives
\[
 e^2\left(\frac{e^2-k^2}{4}-\operatorname{Re}z\right)
    -(\operatorname{Im}z)^2
   =(a-e^2/2)^2+e^2r^2-a^2-b^2\ge0.
\]
We have proved the actual enclosure
\[
 \boxed{\operatorname{Re}z\le\frac{E_N(t)^2-k^2}{4},\qquad
 (\operatorname{Im}z)^2\le E_N(t)^2
       \left(\frac{E_N(t)^2-k^2}{4}-\operatorname{Re}z\right).}
 \tag{LC11}
\]
It includes every eigenvalue by taking its nonzero eigenvector;
no diagonalizability is needed. It includes the case \(E_N(t)=0\).
Every quantity in (LC11) is supplied by the original finite metric,
the monic coefficient functional and the actual deformation parameter.

\subsection{The complete period transport}

For fixed \(u\ne0\), the full determinant proof makes the original
period matrix \(\Pi(u,t)\) invertible at every parameter. Put
\[
 \widetilde G=\Pi^{-*}G\Pi^{-1},\quad
 \widetilde A_t=\Pi A_t\Pi^{-1},\quad
 \widetilde L_t=\Pi L_t\Pi^{-1},\quad
 \widetilde W_t=\Pi^{-*}W_t\Pi^{-1}.
 \tag{LC12}
\]
Direct multiplication gives
\(\widetilde W_t=\widetilde A_t^*\widetilde G+
\widetilde G\widetilde A_t-k\widetilde G\) and
\(\widetilde L_t=\widetilde A_t^2-k\widetilde A_t\).
For \(x=\Pi v\), its norm and the complete numerical-range quotient
are exactly those in (LC10). Its control quotient and adjoint-defect
identity are likewise transported by the literal congruences in
(LC12). Thus (LC11) applies unchanged in that period realization.

The transported fixed action is \(\Pi A\Pi^{-1}\); the displayed
deformed action differs from it by \(t\Pi R\Pi^{-1}\), whose
control contribution was computed in (LC4)--(LC7). The arithmetic
source realizes precisely this change through
\(D_t^{(k)}=D^{(k)}+t r_NRj_E\), with
\(j_Er_N=I\) and the original boundary primitive retained, as
proved in (SC7)--(SC8). These equalities supply the source,
coefficient and period maps for the same Laplacian control.
No smallness uniform in packet size is claimed for the fully
calculated \(\epsilon_N\) or \(\sigma_N\).
